\documentclass[11pt]{article}
\usepackage[a4paper,margin=1in]{geometry}
\usepackage{amsmath,amssymb,amsthm}
\usepackage{hyperref}
\usepackage{enumitem}
\usepackage{parskip}

\title{Tour --- Reinforcing Recursive Language Models\\
       (learning-map companion document)}
\author{Pablo Leyva}
\date{May 15, 2026}

\begin{document}
\maketitle

A guided walk through the learning-map for the alphaXiv blog
\emph{Reinforcing Recursive Language Models}. Read the paper graph plus
improvements graph plus foundation prereqs in the order given here.

\section{Reader's contract}

\textbf{Who this is for.} A reader at the level of an upper-undergraduate
machine learning student --- comfortable with PyTorch, has seen REINFORCE / PPO
once, has read a transformer-pre-training tutorial. The ``advanced'' concept
pages assume the reader can manipulate score-function estimators and clipped
surrogates without consulting a textbook.

\textbf{Expected reading time.} 3--4 hours for the paper graph plus 1--2 hours
for the improvements graph plus 30 minutes for the proofs --- roughly one full
day of focused study.

\textbf{Math-foundations entry point.} If the reader is not yet comfortable
with REINFORCE / PPO / GRPO at the equation level, start at
\texttt{first-principles-to-llms} Chapter~31 and work backwards. If the reader
has not yet built a transformer LM end-to-end, start at Chapter~27. If the MDP
framing is unfamiliar, start at Chapter~29.

\textbf{Pacing convention.} Each prereq is tagged \texttt{skim} (5--10 min
refresh), \texttt{read} (30--45 min careful read), or \texttt{drill} (1--2 h
with paper-and-pencil derivations).

\section{Foundations walk (chain prereqs)}

Topologically-ordered foundation prereqs from the chain repository
(\texttt{pleyva2004/first-principles-to-llms}):
\begin{enumerate}
\item \textbf{Ch 25 --- Causal LM as MDP} (\texttt{skim}). Every node in the RLM
      tree is one Ch 25 MDP; nothing in the paper changes that, only the
      inter-node structure.
\item \textbf{Ch 27 --- Tiny GPT pre-training} (\texttt{skim}). The policy
      $\pi_\theta$ is just an autoregressive transformer; the SFT phase reuses
      Ch 27's NTP loss.
\item \textbf{Ch 28 --- SFT, RLHF/PPO, DPO} (\texttt{read}). The cold-start SFT
      of concept~14 is a direct lift, and the per-node loss of concept~7 is
      PPO. If the reader has not done PPO end-to-end before, this is the
      chapter to spend time on.
\item \textbf{Ch 29 --- MDP foundations} (\texttt{read}). The trajectory tree
      of concept~5 is a tree of MDP rollouts; the Bellman setup is what gets
      lifted to SMDPs in improvement~104.
\item \textbf{Ch 30 --- Max-ent + soft Bellman} (\texttt{skim}). Optional ---
      gives a principled way to add an entropy bonus to the per-node loss; not
      needed for understanding the paper, but useful for extensions.
\item \textbf{Ch 31 --- Policy gradient, GRPO, RLHF/DPO bridge} (\texttt{drill}).
      This is \emph{the} chapter the paper extends. The
      baseline-doesn't-bias proof is what makes inheritance unbiased; the GRPO
      sequence-level credit-assignment trick is what gets lifted to trees. If
      the reader does nothing else from the chain, drill this.
\end{enumerate}

\section{Paper concepts walk}

Topologically-ordered list of the 14 paper concepts (see
\texttt{paper/README.md} for the DAG):
\begin{enumerate}
\item Concept 1 --- Recursive Language Model: what an RLM rollout \emph{is}.
\item Concept 2 --- Python REPL environment: the REPL that hosts every node.
\item Concept 3 --- Root rollout: the only node that gets reward.
\item Concept 4 --- Child rollout (\texttt{rlm\_query}): what the recursion
      primitive produces.
\item Concept 5 --- RLM trajectory tree: the full session is a tree.
\item Concept 6 --- Shared policy: one $\theta$ for both roles.
\item Concept 7 --- PPO clipped surrogate: the per-node atom.
\item Concept 8 --- GRPO advantage: group-relative baseline.
\item Concept 9 --- Advantage inheritance: child's $A := A_g$.
\item Concept 10 --- $1/k_g$ averaging: balancing trick.
\item Concept 11 --- Recursive subtree loss: depth-agnostic loss.
\item Concept 12 --- Evidence selection task: the workload.
\item Concept 13 --- Rubric LLM-judge reward: the training signal.
\item Concept 14 --- Cold-start SFT: bootstrap before RL.
\end{enumerate}

\textbf{Reading strategy.} Concepts 1--6 are scaffolding; spend 10--15 minutes
each. Concepts 7--11 are the load-bearing math; spend 30--45 minutes each, run
the corresponding \texttt{code/} files, and verify the printed output matches
the reader's mental model. Concepts 12--14 are the experimental setup; skim
unless reproducing.

\section{Improvements walk}

Four improvement proposals --- two PROOF-validated math results and two
MEASUREMENT-validated experiments. The DAG is in \texttt{improvements/README.md}.

\subsection*{101. Tighter unbiasedness theorem for advantage inheritance --- PROOF}
The blog asserts inheritance is ``unbiased in the same sense as GRPO'' but
does not formalise the assumption. The proposed theorem makes the
conditional-independence assumption $(*)$ precise, proves unbiasedness under
it, and quantifies bias when it fails (e.g., tool-use chains where child A's
output is child B's input). Validation: \texttt{proofs/inheritance-unbiasedness.tex}.

\subsection*{102. Per-child local-baseline variance reduction --- MEASUREMENT}
Subtract a per-state baseline $b(s_v)$ from the inherited advantage before
applying the surrogate. Standard control-variate trick; under $(*)$ it is
exactly unbiased and reduces variance whenever the baseline correlates with
the score-function-weighted advantage. Validation:
\texttt{improvements/local-baseline.py}.

\subsection*{103. $k_g$ scale-law experiment --- MEASUREMENT}
Sweep $k_g \in \{1, 2, 4, 8, 16, 32\}$ at fixed compute budget; expect an
interior optimum where per-parent reward signal and per-step batch diversity
are balanced. Validation: \texttt{improvements/k-scale-sweep.py}.

\subsection*{104. RLM children as options (Sutton-Precup-Singh) --- PROOF}
Map \texttt{rlm\_query} calls to options $o = (I, \pi_o, \beta)$, lift the
whole RLM training problem to a semi-Markov decision process, and import SMDP
Q-learning convergence (Sutton-Precup-Singh 1999, Theorem~1) under the
natural conditions. The mapping is mostly a translation. Validation:
\texttt{proofs/rlm-as-options.tex}.

\section{What to do next}

\begin{enumerate}
\item \textbf{Reproduce the toy bandit reward-variance gain.} Run
      \texttt{python3 improvements/local-baseline.py}; verify the printed
      ratio is below 1.0. If not, the local baseline is mis-specified.
\item \textbf{Run the $k_g$ sweep on a non-toy task.} The bandit
      interior-optimum is a small effect. Re-run the sweep on a real
      arxiv-paper retrieval workload (or one of the reader's own) to see if
      the optimum is at $k = 4$, $k = 16$, or larger. The answer determines
      deployment cost.
\item \textbf{Read \texttt{proofs/rlm-as-options.tex} and identify the most
      promising corollary to prove next.} The SMDP framing imports a lot of
      theory for free; the question is which corollary (interruptibility,
      intra-option learning, ...) is most actionable for the next paper.
      Pick one and write a 5-page extension.
\end{enumerate}

\end{document}
